The method requires only a pre-trained flow-based model and a numerical ODE/SDE solver. The three samplers of Section~\ref{subsec:supplying_diffusion} share a single loop and differ only in the diffusion $\gdiff_\tau$, the guidance weight $\gweight_\tau=\vscoef_\tau+\tfrac12\gdiff_\tau^2$, the source, and whether a Brownian increment is added. Algorithm~\ref{alg:unified} states the shared loop for one assimilation step, i.e.\ drawing a posterior sample $\state^{n}\sim\conddist{\state^{n}}{\obs^{n},\state^{n-1}}$. Algorithms~\ref{alg:si_sde}--\ref{alg:fm_ode} specialise it to the three samplers; the differences are confined to the boxed lines.

\begin{algorithm}[ht]
\caption{Unified posterior assimilation step (shared loop for all three samplers)}
\label{alg:unified}
\begin{algorithmic}[1]
\STATE \textbf{Input:} previous state $\bx_0=\state^{n-1}$; observation $\obs=\obs^{n}$; trained model (SI drift $\drift_\theta$ or FM velocity $\fmvel_\theta$); schedules $\alpha_\tau,\beta_\tau,\sigpath_\tau$; diffusion $\gdiff_\tau$; steps $M$
\STATE \textbf{Sampler-specific:} set source/anchor, the diffusion $\gdiff_\tau$, and weight $\gweight_\tau=\vscoef_\tau+\tfrac12\gdiff_\tau^2$ as in Alg.~\ref{alg:si_sde}, \ref{alg:fm_sde}, or \ref{alg:fm_ode}
\STATE Initialise $\bXstar$ from the (reweighted) source; set $\Delta\tau = 1/M$
\FOR{$\tau = 0,\ \Delta\tau,\ \ldots,\ 1-\Delta\tau$}
    \STATE Evaluate velocity $\fmvel_\tau$ and score $\genscore_\tau$ from the model \hfill // Eq.~\eqref{eq:SI_score} (SI) or \eqref{eq:fm_score} (FM)
    \STATE Form prior drift $\drift^{\gdiff}_\tau = \fmvel_\tau + \tfrac12\gdiff_\tau^2\genscore_\tau$ \hfill // $=\drift_\theta$ for SI with $\gdiff_\tau=\gamma_\tau$
    \STATE Compute observation interpolant $\interpolantobs_\tau = \alpha_\tau\obsmatrix\bm a_0 + \beta_\tau\obs$
    \STATE Source moments $\srcmean_\tau,\srccov_\tau$; likelihood mean $\bar\mu_\tau=\obsmatrix\bXstar-\obsmatrix\srcmean_\tau$ \hfill // Lemma~\ref{lem:source_moments}
    \STATE Interpolant score $\bar S_\tau = (\nabla_{\bx}\bar\mu_\tau)^\top\bar\Sigma_\tau^{-1}(\interpolantobs_\tau-\bar\mu_\tau)$; gain $G_\tau$ \hfill // Eq.~\eqref{eq:multiplicative_correction}; $\nabla_{\bx}\bar\mu_\tau\approx\obsmatrix$ Jacobian-free
    \STATE Posterior drift $\postdrift_\tau = \drift^{\gdiff}_\tau + \boxed{\gweight_\tau}\, G_\tau\bar S_\tau$
    \STATE Update $\bXstar \leftarrow \bXstar + \postdrift_\tau\,\Delta\tau + \boxed{\gdiff_\tau\sqrt{\Delta\tau}\,\bz},\ \ \bz\sim\mathcal{N}(0,\bI)$ \hfill // noise term omitted when $\gdiff_\tau=0$
\ENDFOR
\STATE \textbf{Return:} $\state^{n} = \bXstar$
\end{algorithmic}
\end{algorithm}

\begin{algorithm}[ht]
\caption{Sampler 1 -- Stochastic interpolant SDE (native diffusion)}
\label{alg:si_sde}
\begin{algorithmic}[1]
\STATE \textbf{Model:} SI drift $\drift_\theta$; \textbf{anchor} $\bm a_0=\bx_0$; \textbf{diffusion} $\gdiff_\tau=\gamma_\tau$ (native)
\STATE \textbf{Weight} $\gweight_\tau = \vscoef_\tau + \tfrac12\gamma_\tau^2 = \gamma_\tau^2\tau\big(\tfrac{\dot\beta_\tau}{\beta_\tau}\gamma_\tau-\dot\gamma_\tau\big)\big/\gamma_\tau$ \hfill // velocity--score coef.\ for $\sigpath_\tau=\gamma_\tau\sqrt\tau$
\STATE \textbf{Init} $\bXstar\leftarrow\bx_0$ \hfill // point-mass source; reweighting vacuous
\STATE \textbf{Drift} $\drift^{\gdiff}_\tau=\drift_\theta(\bXstar,\bx_0,\tau)$ \hfill // score already built in
\STATE \textbf{Source moments} via Wiener identities (Lemma~\ref{lem:source_moments}, SI case)
\STATE Run Algorithm~\ref{alg:unified} with the Brownian increment \emph{retained}.
\end{algorithmic}
\end{algorithm}

\begin{algorithm}[ht]
\caption{Sampler 2 -- Flow matching SDE (lifted diffusion)}
\label{alg:fm_sde}
\begin{algorithmic}[1]
\STATE \textbf{Model:} FM velocity $\fmvel_\theta$; \textbf{anchor} $\bm a_0=\bm 0$, $\bx_0$ as conditioning input; \textbf{diffusion} $\gdiff_\tau>0$ (free, e.g.\ $\gdiff_\tau\propto\sqrt{\alpha_\tau\beta_\tau}$)
\STATE \textbf{Weight} $\gweight_\tau = \vscoef_\tau + \tfrac12\gdiff_\tau^2$, with $\vscoef_\tau=\alpha_\tau(\dot\beta_\tau\alpha_\tau-\dot\alpha_\tau\beta_\tau)/\beta_\tau$
\STATE \textbf{Init} $\bXstar\sim\mathcal{N}(0,\bI)$ \hfill // exact for marginal construction ($\Phi_0^{\obs}$ const)
\STATE \textbf{Score} $\genscore_\tau$ recovered from $\fmvel_\theta$ via Eq.~\eqref{eq:fm_score}; drift $\drift^{\gdiff}_\tau=\fmvel_\theta+\tfrac12\gdiff_\tau^2\genscore_\tau$
\STATE \textbf{Source moments} via Tweedie identities (Lemma~\ref{lem:source_moments}, FM case)
\STATE Run Algorithm~\ref{alg:unified} with the Brownian increment \emph{retained}.
\end{algorithmic}
\end{algorithm}

\begin{algorithm}[ht]
\caption{Sampler 3 -- Guided probability-flow ODE (flow matching, deterministic)}
\label{alg:fm_ode}
\begin{algorithmic}[1]
\STATE \textbf{Model:} FM velocity $\fmvel_\theta$; \textbf{anchor} $\bm a_0=\bm 0$; \textbf{diffusion} $\gdiff_\tau=0$
\STATE \textbf{Weight} $\gweight_\tau = \vscoef_\tau = \alpha_\tau(\dot\beta_\tau\alpha_\tau-\dot\alpha_\tau\beta_\tau)/\beta_\tau$ \hfill // pure velocity--score term
\STATE \textbf{Init} $\bXstar\sim\mathcal{N}(0,\bI)$ \hfill // randomness enters only through the source
\STATE \textbf{Drift} $\drift^{\gdiff}_\tau=\fmvel_\theta(\bXstar,\state^{n-1},\tau)$ \hfill // no score term needed ($\gdiff_\tau=0$)
\STATE \textbf{Source moments} via Tweedie identities (Lemma~\ref{lem:source_moments}, FM case)
\STATE Run Algorithm~\ref{alg:unified} with the Brownian increment \emph{omitted} (deterministic Euler step / higher-order ODE solver).
\end{algorithmic}
\end{algorithm}

Several points are worth noting. First, the loop never solves an inner integral: the likelihood score is read off in closed form from the observation interpolant, and the gain $G_\tau$ is a rank-$N_y$ correction. With the Jacobian-free choice of Corollary~\ref{cor:cheap_drift}, $\obsmatrix^\top\obscov^{-1}\obsmatrix$ is precomputed once and each step costs one model evaluation plus an $\mathcal{O}(N_u N_y)$ correction, identically for all three samplers. Second, the dynamics are continuous in pseudo-time, so the solver and number of steps $M$ may be chosen after training to trade accuracy against cost; the guided ODE admits higher-order deterministic solvers (e.g.\ Heun, RK4), while the SDE samplers use Euler--Maruyama or stochastic Runge--Kutta. Third, the interpolant coefficients degenerate at $\tau=0$: since $\beta_0=0$, the gain $G_\tau$ and the weight $\gweight_\tau$ are singular there. In practice the correction is applied from $\tau=\Delta\tau$ onwards. Fourth, an ensemble of posterior samples is obtained by running the loop with independent draws (independent Brownian paths for the SDE samplers, independent initial latents for all FM samplers); a full trajectory follows by feeding each posterior sample back as $\state^{n}$ and assimilating $\obs^{n+1}$. Finally, the two FM samplers initialise from $\mathcal{N}(0,\bI)$, which is exact for the marginal construction: at $\tau=0$ the latent is independent of $\bx_1$, so $\Phi_0^{\obs}$ is constant and no reweighting of the initial latents is needed (Section~\ref{subsec:supplying_diffusion}, Appendix~\ref{appendix:proof_conditioning}). The guided ODE relies entirely on this source randomness, since its dynamics inject none.
